% Generated draft: requires mathematical and source review.
\documentclass[11pt]{article}
\usepackage[margin=25mm]{geometry}
\usepackage{fontspec}
\setmainfont{Latin Modern Roman}
\usepackage{amsmath,amssymb,mathtools,unicode-math}
\setmathfont{Latin Modern Math}
\usepackage{enumitem,needspace,xurl,hyperref}
\setlength{\parindent}{0pt}
\setlength{\parskip}{5pt}
\setlength{\emergencystretch}{3em}
\newcommand{\N}{\mathbb{N}}
\newcommand{\Z}{\mathbb{Z}}
\newcommand{\Q}{\mathbb{Q}}
\newcommand{\R}{\mathbb{R}}
\newcommand{\C}{\mathbb{C}}
\newcommand{\F}{\mathbb{F}}
\newcommand{\A}{\mathbb{A}}
\newcommand{\PP}{\mathbb{P}}
\newcommand{\E}{\mathbb{E}}
\newcommand{\eps}{\varepsilon}
\newcommand{\dd}{\,\mathrm{d}}
\begin{document}
\section*{UnsolvedMath ID: 701499}
\textbf{Source problem code:} catalogue entry
\subsection{Definitions and mathematical statement}
A group \(W\) is a Coxeter group if it admits a presentation
\[
W=\left\langle S\ \middle|\ s^2=1\ (s\in S),\ (st)^{m_{st}}=1\ (s\ne t,\ m_{st}<\infty)\right\rangle,
\]
where \(m_{st}=m_{ts}\in\{2,3,\ldots\}\cup\{\infty\}\) for distinct \(s,t\in S\). The value \(m_{st}=\infty\) imposes no relation on \(st\). A group is finitely generated if some finite subset generates it.

A group \(G\) is conjugacy separable if, whenever \(x,y\in G\) are not conjugate in \(G\), there exist a finite group \(Q\) and a surjective homomorphism \(\pi:G\to Q\) such that \(\pi(x)\) and \(\pi(y)\) are not conjugate in \(Q\). Thus the required separation condition is
\[
\left(\nexists g\in G:\ y=gxg^{-1}\right)
\Longrightarrow
\left(\exists Q\text{ finite},\ \exists\pi:G\twoheadrightarrow Q:\
\nexists q\in Q:\ \pi(y)=q\pi(x)q^{-1}\right).
\]

\textbf{Problem (A. Minasyan, 2014).} Is every finitely generated Coxeter group conjugacy separable?

The supplied record labels this problem open and states that some special cases were considered in the paper by P.-E. Caprace and A. Minasyan cited below. It does not specify those special cases. These source qualifications are retained without an independent assessment of the literature.
\subsection{Short English statement}
Must every pair of nonconjugate elements of a finitely generated Coxeter group remain nonconjugate in some finite quotient? This is the question attributed to A. Minasyan; its open status is supplied by the registry.
\subsection{Sources}
\begin{itemize}
\item Local UnsolvedMath registry, supplied source record: permanent ID \(701499\), exact source code \texttt{catalogue entry}, title \emph{Kourovka Notebook Problem 18.70}. The record labels the problem open and attributes it to A. Minasyan (2014).
\item \emph{The Kourovka Notebook (Unsolved Problems in Group Theory)}, original problem 18.70, 18th issue (2014); the supplied background also identifies the 21st issue edition (2026), arXiv:1401.0300v48. The supplied collection description names E. I. Khukhro and V. D. Mazurov as editors.
\item P.-E. Caprace and A. Minasyan, \emph{Illinois J. Math.} \textbf{57}, no. 2 (2013), 499--523. The supplied statement says that some special cases were considered there; it supplies neither an article title nor a description of those cases.
\end{itemize}
\end{document}
